\subsection{The Auxiliary Persistence Layer (Redis)}
\label{sec:redis_architecture}

LangGraph's built-in checkpointing writes a snapshot of the graph state after every node transition, which is enough for fault tolerance within a single session. It does not, however, retain user preferences or project context across separate sessions: once a thread terminates, all accumulated knowledge, such as the last used MCU board or the preferred compression level, is lost.

To fill this gap, I introduced an auxiliary persistence layer built on Redis. It operates on two levels. The first, \textit{short-term memory}, is the session state managed by LangGraph's \texttt{AsyncRedisSaver}. It stores the transient graph snapshot that enables HITL interrupt/resume and time-travel debugging. The second, \textit{long-term memory}, is a custom storage layer keyed by \texttt{user\_id}. It survives across sessions and retains context such as preferred hardware boards, last used models, and past project paths.

Together these two layers enable \textit{context chaining}: outputs from a previous session (e.g., the MCU target selected during \texttt{firmware\_flow}) automatically pre-populate the input context of future workflows (e.g., \texttt{ai\_flow}), eliminating redundant data entry.

\subsubsection{Short-Term Memory (LangGraph Checkpoints)}

LangGraph checkpoints are stored using a composite key built from the \texttt{user\_id}, the \texttt{session\_id} (identifying the client instance), and a timestamp-based \texttt{checkpoint\_id} managed automatically by LangGraph, as in the following example:

\begin{verbatim}
checkpoint:michele:vscode-session:1707575432123
checkpoint:michele:vscode-session:1707575455789
checkpoint:john:cli-session:1707575500000
\end{verbatim}

This hierarchical naming enforces strict user isolation: different users can never access each other's checkpoints. It also supports session multiplexing, so a single user can run independent workflows from VS Code and CLI at the same time. Because each checkpoint ID contains a timestamp, the system preserves a clear temporal ordering that enables time-travel debugging.

\subsubsection{Long-Term Memory (User Profiles)}

User profiles follow a simpler key structure for direct access (e.g., \texttt{user:michele:profile}). Each profile is a JSON-serialized dictionary that captures the user's most recent working context, as shown in Listing~\ref{lst:profile_schema}.

\begin{lstlisting}[language=json, caption={User profile schema stored in Redis}, label={lst:profile_schema}]
{
  "board_name": "NUCLEO-H743ZI",
  "mcu_series": "STM32H7",
  "last_model": "mobilenet_v1_0.25_128",
  "last_project_path": "/home/user/workspace/h7_project",
  "last_workflow": "firmware",
  "timestamp": "2024-02-10T14:35:22Z"
}
\end{lstlisting}

The schema records the most recently used STM32 board and MCU series (for pre-filling firmware generation requests), the last AI model selected for optimization, and the path of the latest firmware project directory. It also stores the last executed workflow route and a timestamp marking the most recent successful completion. When a new session starts, these fields allow the orchestrator to skip redundant questions and resume from a familiar context.

\subsubsection{Redis Integration and Implementation}

The persistence layer is initialized during FastAPI server startup. An asynchronous Redis client connects to the local \texttt{redis://localhost:6379} instance and is wrapped by LangGraph's \texttt{AsyncRedisSaver}, which handles all checkpoint I/O transparently:

\begin{lstlisting}[language=Python, caption={Redis checkpointer initialization in server.py}, label={lst:redis_init}]
from langgraph.checkpoint.redis.aio import AsyncRedisSaver
import redis.asyncio as redis

# 1. Create async Redis client
checkpointer_redis = redis.from_url(
    "redis://localhost:6379/0",
    encoding="utf-8",
    decode_responses=True
)

# 2. Initialize AsyncRedisSaver
memory = AsyncRedisSaver(redis_client=checkpointer_redis)

# 3. Setup internal RedisVL indexes (run once on startup)
await memory.setup()

# 4. Compile graph with checkpointer
graph = builder.compile(checkpointer=memory)
\end{lstlisting}

The \texttt{memory.setup()} call creates internal Redis indexes used by RedisVL for efficient checkpoint retrieval and time-travel queries. Each time a graph node completes, LangGraph automatically persists the \texttt{MasterState} to Redis as a MessagePack or JSON-serialized payload:

\begin{verbatim}
{
  "v": 1,
  "ts": "2024-02-10T15:30:45.123Z",
  "channel_values": {
    "messages": [...],
    "firmware_project_dir": "/tmp/stm32_project",
    "model_path": "/cache/mobilenet_v2.h5",
    "route": "ai_flow",
    "nni_search_space": {...}
  },
  "parent_config": {
    "configurable": {
      "thread_id": "michele:vscode-session",
      "checkpoint_id": "1707575432123"
    }
  }
}
\end{verbatim}

The \texttt{channel\_values} field contains the complete \texttt{MasterState} with all 50+ attributes, \texttt{parent\_config} holds the reference to the previous checkpoint (for time-travel), and \texttt{ts} provides the timestamp for ordering.

The long-term memory layer is managed directly in \texttt{server.py}. Before graph execution, the server loads the user profile from Redis and injects it into the initial state:

\begin{lstlisting}[language=Python, caption={User profile loading in server.py}, label={lst:profile_load}]
# 1. Construct profile key
user_profile_key = f"user:{request.user_id}:profile"

# 2. Fetch from Redis
raw_profile = await redis_client.get(user_profile_key)
user_profile = json.loads(raw_profile) if raw_profile else {}

# 3. Inject into initial state
initial_state = {
    "message": last_user_message,
    "persistent_context": user_profile,
    "reset_profile": False
}
\end{lstlisting}

The \texttt{persistent\_context} field is then accessible to all graph nodes for context-aware decisions. After workflow completion or critical state mutations, the profile is updated with the latest values and written back:

\begin{lstlisting}[language=Python, caption={User profile saving in server.py}, label={lst:profile_save}]
# 1. Extract current values from state
new_profile = {
    "board_name": state.get("board_name"),
    "mcu_series": state.get("mcu_series"),
    "last_model": state.get("selected_model"),
    "last_project_path": state.get("firmware_project_path") or state.get("firmware_project_dir"),
    "last_workflow": state.get("route"),
    "timestamp": state.get("timestamp")
}

# 2. Update profile dictionary
new_profile = {k: v for k, v in new_profile.items() if v is not None and v != ""}
updated_profile = {**user_profile, **new_profile}

# 3. Serialize and save to Redis
await redis_client.set(user_profile_key, json.dumps(updated_profile))
\end{lstlisting}

\subsubsection{System Interaction Scenarios}

The following scenarios illustrate how the dual-memory architecture supports real-world usage patterns.

\paragraph{Scenario 1: Multi-Day NNI Experiment}

\textbf{Day 1, Session 1 (VS Code):}\newline
The user requests an NNI experiment on MobileNetV2 for the Fruit-360 dataset. The framework launches 20 NNI trials, estimated to take 12 hours. After 2 hours the user closes VS Code. The short-term memory saves the checkpoint with \texttt{route = "nni\_optimization"} and \texttt{nni\_status = "running"}, while the long-term memory records \texttt{last\_board = "STM32H7"}.

\textbf{Day 2, Session 2 (CLI):}\newline
The user reconnects via CLI using the same \texttt{user\_id} and \texttt{session\_id}. The server calls \texttt{graph.aget\_state(config)}, retrieves the Day 1 checkpoint, and detects \texttt{current\_\allowbreak state.next = ["check\_\allowbreak nni\_\allowbreak status"]}. The graph resumes automatically: NNI results are retrieved from disk, the best model is saved, and integration proceeds.

\paragraph{Scenario 2: Board Preference Pre-Fill}

\textbf{Session A (Firmware Generation):}\newline
The user creates a project for STM32H743. Once firmware is generated, the profile is updated with \texttt{last\_board = "STM32H743ZI"}.

\textbf{Session B (AI Model Selection, one week later):}\newline
The user requests ``Analyze MobileNetV2'' without specifying hardware. The \texttt{collect\_\allowbreak analysis\_\allowbreak info()} node reads \texttt{state.last\_board} from the persistent context and auto-populates \texttt{target\_mcu = "stm32h743"}, eliminating one clarification roundtrip.

\paragraph{Scenario 3: HITL Interrupt and Resume}

During execution, a workflow node calls \texttt{interrupt(\{type: "decision\_request", message: "Proceed to AI?"\})}. LangGraph suspends and saves a checkpoint with \texttt{next = ["decide\_\allowbreak continue\_\allowbreak to\_\allowbreak ai"]}. The API server forwards a JSON event to VS Code, where the user sees: ``Firmware generated. Continue with AI model selection?''

When the user replies (e.g., ``Yes, use high compression''), VS Code sends a \texttt{POST /stream} request. The server retrieves the latest checkpoint, injects the user response via \texttt{graph.aupdate\_state()}, and the graph resumes from the exact node that issued the interrupt, parsing the reply through an LLM and routing accordingly.

\subsubsection{Performance Considerations}

\paragraph{Redis I/O Overhead}

Checkpoint writes inherently add latency to the execution pipeline, typically in the range of 5 to 20 milliseconds per node transition. To mitigate this overhead, the framework employs several strategies, such as the use of non-blocking asynchronous I/O operations in the \texttt{AsyncRedisSaver} to prevent main event loop stalls. Network transit latency is minimized during development by defaulting to a local \texttt{redis://localhost} connection. LangGraph further optimizes this process by consolidating multiple internal state updates from a single node's execution context into a single write operation before transitioning to the next stage of the graph.

\paragraph{Memory Footprint and Evolution of Data Persistence}

During early development, the system attempted to store large artifacts (such as generated firmware zips, serialized models, and base64 images) directly within the Redis state. However, this approach presented technical bottlenecks, as Redis string size limits and synchronization overhead slowed down the orchestration loop.

To resolve this, the architectural design evolved toward a hybrid persistence model using \textbf{Container Persistent Volumes}. Instead of using an intermediate database for heavy files, the system writes artifacts directly to persistent storage volumes mapped to individual container instances. This strategy perfectly aligns with a Cloud-Native architecture: in a future Kubernetes deployment, each user will operate within a dedicated namespace with dynamically provisioned PersistentVolumeClaims (PVCs).

Redis is now strictly limited to maintaining the lightweight conversational checkpointing of the LangGraph execution thread, while the storage of project files is handled natively by the file system. This simplifies the architecture, improves I/O performance, and guarantees environment isolation without relying on database abstractions.

\subsubsection{Security Considerations}

Given that user profiles may store sensitive information, such as STMicroelectronics developer credentials required for STEdgeAI cloud authentication, the framework must adhere to strict security practices. In production deployments, it is imperative to encrypt sensitive fields using solid algorithms like AES-256 before persisting them in the Redis data store.

Also, access control should be strictly enforced using Redis ACL (Access Control Lists) to restrict client permissions, ensuring that specific key patterns (e.g., \texttt{user:*:profile}) are only accessible to authenticated and authorized services.

Finally, to comply with GDPR data minimization principles and prevent indefinite data retention, user profiles are configured with a Time-To-Live (TTL) expiration policy, ensuring they are automatically purged after one year of complete inactivity.
